\section{Limitations}
\label{sec:limitations}

Stated plainly, because a service whose limits are hidden is worse than one
whose limits are known.

\begin{enumerate}[leftmargin=1.8em,itemsep=0.7ex]

\item \textbf{The footprints are not Bhoonidhi's.} Its \ep{/data/search} needs a
JWT and no credentials were available, so footprints come from Microsoft
Planetary Computer's public STAC --- \textbf{68 real acquisition footprints},
60 optical and 8 radar, with real timestamps and cloud cover. Nothing is
synthesised. This is labelled, not hidden: the store carries
\code{is\_substitute\_for\_bhoonidhi: true} and that flag propagates into
\emph{every process response}. Supplying credentials and re-running
\code{fetch\_footprints.py --source bhoonidhi} changes nothing else --- same
code path, same processes, same tests. \textbf{This is the one thing to fix if
an account becomes available.} The raster is substituted the same way
(Copernicus DEM GLO-30 standing in for CartoDEM).

\item \textbf{Cloud cover is reported but not used.} \code{aoi-coverage} counts
a footprint as coverage however cloudy it is, and \code{minimal-scene-set} will
select a 90\,\%-cloud scene if it contributes most area. So the 92\,\% figure
is about \emph{geometric} redundancy only. Cloud-aware coverage is the obvious
and most valuable next step.

\item \textbf{Greedy set cover is not guaranteed minimal} --- within a $\ln(n)$
factor of optimal, and the response says so in its \code{algorithm} field.

\item \textbf{The catalogue tier is collection-level} (59 collections, not the
granules behind them) and its extents are \textbf{indicative and labelled as
such}: temporal start is the mission launch date, a defensible lower bound since
data cannot predate launch. Every record carries
\code{spatial\_extent\_basis} and \code{temporal\_start\_basis}. Good to
discover with, not to cite.

\item \textbf{No authentication, quota or rate limiting.} The processes are
bounded (caps on footprints matched and grid cells built) but there is no
per-caller quota. Fine for a graded demo on an ephemeral URL; not for a public
service, which would need an API gateway.

\item \textbf{No CQL2 on the catalogue tier} --- pygeoapi's TinyDB provider
gives full-text, property equality, bbox, datetime and sorting. PostgreSQL
would add the Filtering class. This limits the supporting tier, not Processes.

\item \textbf{One raster over one window}, and \textbf{no time-series or change
detection}: processes filter footprints by a \code{start}/\code{end} window but
nothing models change \emph{between} windows.

\item \textbf{The public URL is ephemeral} --- a Cloudflare quick tunnel gets a
fresh hostname each run. A stable URL needs a named tunnel or a cloud VM. The
tested deployment is \ep{http://localhost:5003}, which the brief permits.

\item \textbf{Self-written compliance suite, not OGC CITE}
(Section~\ref{sec:compliance}), and \textbf{not an official NRSC product}. This
mirrors no data: footprint geometry and metadata only, never pixels from
Bhoonidhi.

\end{enumerate}

\paragraph{What would be built next, in order.} Cloud-aware coverage, so the
service optimises for \emph{usable} rather than merely intersecting imagery
(this would change the results, not just extend them); an authenticated
Bhoonidhi harvest, removing the substitution; granule-level records with CQL2
over PostgreSQL; an exact set-cover solver for small candidate sets, falling
back to greedy; and a quota gateway ahead of any genuinely public deployment.
